\chapter{Systemarchitektur}
\label{chap:systemarchitektur}

Basierend auf den im vorherigen Kapitel ermittelten Anforderungen wird im Folgenden die 
Systemarchitektur des drahtlosen Beleuchtungssystems vorgestellt. 
Die Architektur beschreibt die grundlegende Struktur des Systems, die ausgewählten 
Hardwarekomponenten sowie deren Zusammenwirken zur Realisierung der angestrebten Funktionalität. 
Während dieses Kapitel die konzeptionellen und strukturellen Entscheidungen – insbesondere 
hinsichtlich der Hardware – aufzeigt, wird die konkrete technische Umsetzung 
in Kapitel~\ref{chap:implementierung} (Implementierung) behandelt.

\section{Systemüberblick}
\label{sec:systemueberblick}

Das entworfene System setzt sich aus zwei Hauptkomponenten zusammen: einer zentralen Sendeeinheit (Bridge) und 
mehreren dezentralen Empfangseinheiten (LED-Leuchtröhren). 
Die Bridge fungiert als zentraler Kommunikationsknotenpunkt, welcher externe Steuersignale 
aus der Veranstaltungstechnik – wahlweise über DMX oder Art-Net – entgegennimmt, verarbeitet und in 
ein systeminternes Protokoll übersetzt. 
Anschließend verteilt sie die Daten drahtlos an die angeschlossenen Leuchtröhren.

Die LED-Leuchtröhren agieren als verteilte Aktoren. 
Sie empfangen die von der Bridge ausgesendeten Steuerdaten und setzen diese unmittelbar in Lichtsignale um. 
Darüber hinaus verfügen sie über einen autarken Standalone-Modus, bei dem keine Verbindung zur Bridge notwendig ist. 

Aus dieser Aufteilung ergeben sich die im Folgenden beschriebenen Strukturen: der
Signalfluss je Betriebsmodus, die Auswahl der Hardwarekomponenten, die Gliederung der
Software sowie die Verfahren, mit denen die Daten zwischen den Komponenten übertragen
und verarbeitet werden.

Abbildung~\ref{fig:systemueberblick} zeigt die Komponenten des Systems und die
Verbindungen zwischen ihnen.

\begin{figure}[H]
    \centering
    \begin{tikzpicture}[node distance=1.1cm and 2.3cm]
        \node[extern] (pult) {Lichtsteuerung\\(Pult oder Software)};
        \node[device, right=of pult] (bridge) {\textbf{Bridge}\\ESP32};
        \node[archbox, right=of bridge, yshift=1.7cm] (t1)
            {Leuchtröhre 1\\\footnotesize XIAO ESP32-S3 + WS2815};
        \node[archbox, right=of bridge] (t2)
            {Leuchtröhre 2\\\footnotesize XIAO ESP32-S3 + WS2815};
        \node[archbox, right=of bridge, yshift=-1.7cm] (tn)
            {Leuchtröhre $n$\\\footnotesize XIAO ESP32-S3 + WS2815};

        \draw[link] (pult) -- node[beschr, above] {DMX512\\Art-Net} (bridge);
        \draw[funk] (bridge) -- (t1);
        \draw[funk] (bridge) -- node[beschr, above] {ESP-NOW} (t2);
        \draw[funk] (bridge) -- (tn);

        \node[beschr, below=0.15cm of bridge] {lokales Bedienmenü};
        \node[beschr, below=0.15cm of tn] {lokales Bedienmenü};
    \end{tikzpicture}
    \caption{Systemüberblick: Die Bridge verteilt das kabelgebunden empfangene Steuersignal per Broadcast an alle Leuchtröhren.}
    \label{fig:systemueberblick}
\end{figure}

\section{Betriebsmodi und Signalfluss}
\label{sec:betriebsmodi-signalfluss}

Der Signalfluss im System variiert je nach Betriebsmodus, wobei die grundlegenden Datenpfade und 
Interaktionen zwischen den Komponenten im Folgenden erläutert werden.

 \subsection{Standalone-Modus (Leuchtröhre)}
\label{subsec:signalfluss-einzelbetrieb}

Im Standalone-Modus agieren die LED-Leuchtröhren autark, ohne eine Verbindung zur Bridge.
Die Steuerung erfolgt über die integrierte Benutzeroberfläche der Röhren, bestehend aus 
einem OLED-Display und Navigationsschaltern.
Der interne Controller der Röhre generiert die Lichtsignale basierend auf den Benutzereingaben 
und schickt diese direkt an die LEDs weiter.
In diesem Modus findet keine Kommunikation mit der Bridge statt, und die Röhren arbeiten unabhängig voneinander.


\subsection{DMX-Modus}
\label{subsec:signalfluss-dmx-betrieb}

Im DMX-Modus empfängt die Bridge die Steuerdaten über die DMX- oder die Art-Net-Schnittstelle, verarbeitet diese und leitet 
sie drahtlos an die Leuchtröhren weiter.
Die Bridge fungiert als DMX-zu-ESP-NOW-Gateway: Die DMX-Daten werden in ein internes, an
DMX angelehntes Protokoll übersetzt und über das integrierte Wi-Fi-Modul an die Röhren
gesendet.
Der Controller der Röhre empfängt die Daten über das integrierte Wi-Fi-Modul, validiert diese und setzt sie 
in die entsprechenden LED-Signale um.
Die Steuerbefehle pro Röhre enthalten 4 Bytes. 1 Byte für die Wahl eines vordefinierten Effekts und 3 Bytes für die RGB-Farbwerte.

Welche der beiden Eingangsschnittstellen die Bridge auswertet, wird an ihrem Bedienmenü als
eigener Eintrag der Moduswahl festgelegt: Der DMX-Modus liest das Universum über die
RS-485-Schnittstelle ein, der Art-Net-Modus über die Ethernet-Schnittstelle. Diese
Unterscheidung betrifft ausschließlich die Empfangsstufe der Bridge. Ab der Überführung in
das interne Transportformat sind beide Pfade identisch. Für die Leuchtröhren ist die
Herkunft des Universums nicht erkennbar, sie verbleiben in beiden Fällen im DMX-Modus.
\begin{figure}[H]
	\centering
	\includegraphics[width=0.4\textwidth]{images/diagramm_dmx_mode.png}
	\caption{Signalfluss im DMX-Modus}
	\label{fig:signalfluss-dmx-betrieb}
\end{figure}

\section{Hardware-Architektur}
\label{sec:hardware-architektur}

Die Auswahl der Hardwarekomponenten erfolgte primär unter den Gesichtspunkten der Rechenleistung für 
Echtzeitverarbeitung, der Verfügbarkeit notwendiger Kommunikationsschnittstellen sowie der räumlichen und 
energetischen Restriktionen. Als Basis für beide Systemkomponenten wurden Mikrocontroller der ESP32-Familie gewählt. 
Diese bieten neben einer hohen Taktrate und ausreichend GPIO-Pins vor allem eine integrierte Wi-Fi-Schnittstelle, welche 
die geforderte drahtlose Vernetzung ohne zusätzliche externe Funkmodule ermöglichen und somit das Systemdesign verschlanken.
Abbildung~\ref{fig:hardware-architektur} stellt den Aufbau beider Geräte gegenüber.

\begin{figure}[H]
    \centering
    \begin{tikzpicture}[x=1cm, y=1cm]
        % --- Bridge ---
        \node[archbox, minimum width=3.2cm] (in)   at (0,0)      {DMX512 (RS-485)\\Art-Net (Ethernet)};
        \node[device,  minimum width=2.8cm] (esp)  at (4.9,0)    {ESP32-32U};
        \node[archbox, minimum width=2.4cm] (ant1) at (9.2,0)    {externe Antenne};
        \node[archbox, minimum width=2.8cm] (ui1)  at (4.9,-1.9) {OLED-Display\\+ 3 Taster};

        \draw[link] (in)  -- (esp);
        \draw[link] (esp) -- (ant1);
        \draw[link, {Latex[length=2mm]}-{Latex[length=2mm]}] (esp) -- (ui1);
        \node[beschr, right=0.25cm of ui1] {Versorgung: 5\,V über USB-C};

        \begin{scope}[on background layer]
            \node[huelle, fit=(in)(esp)(ant1)(ui1),
                  label={[font=\small\bfseries]above:Bridge}] {};
        \end{scope}

        % --- Leuchtröhre ---
        \node[archbox, minimum width=3.2cm] (strip) at (0,-5.2)   {WS2815-Strip\\\footnotesize 120 Pixel};
        \node[device,  minimum width=2.8cm] (xiao)  at (4.9,-5.2) {XIAO ESP32-S3};
        \node[archbox, minimum width=2.4cm] (ant2)  at (9.2,-5.2) {externe Antenne};
        \node[extern,  minimum width=3.2cm] (pwr2)  at (0,-7.3)   {12\,V-Netzteil};
        \node[archbox, minimum width=2.8cm] (buck)  at (4.9,-7.3) {Step-Down\\\footnotesize 12\,V $\rightarrow$ 5\,V};
        \node[archbox, minimum width=2.4cm] (ui2)   at (9.2,-7.3) {OLED-Display\\+ 3 Taster};

        \draw[link] (xiao) -- (strip);
        \draw[link] (ant2) -- (xiao);
        \draw[link] (pwr2) -- (strip);
        \draw[link] (pwr2) -- (buck);
        \draw[link] (buck) -- (xiao);
        \draw[link, {Latex[length=2mm]}-{Latex[length=2mm]}] (ui2) -- (xiao);

        \begin{scope}[on background layer]
            \node[huelle, fit=(strip)(xiao)(ant2)(pwr2)(buck)(ui2),
                  label={[font=\small\bfseries]above:Leuchtröhre}] {};
        \end{scope}

        \draw[funk] (ant1) -- node[beschr, right=1mm, pos=0.62] {ESP-NOW} (ant2);
    \end{tikzpicture}
    \caption{Hardware-Architektur beider Gerätetypen: gemeinsame Bedienhardware und gemeinsamer Funkweg, unterschiedliche Schnittstellen.}
    \label{fig:hardware-architektur}
\end{figure}

\subsection{Bridge}
\label{subsec:bridge-hardware}

Die Bridge übernimmt die Aufgabe der Protokollkonvertierung und Datenverteilung. 
Als zentrale Steuereinheit kommt ein ESP32-32U Entwicklungsboard zum Einsatz. Dieser Mikrocontroller bietet 
die entsprechende Rechenleistung um die Echtzeitverarbeitung der eingehenden DMX- und Art-Net-Datenströme zu gewährleisten, 
und verfügt über ausreichend GPIOs für die Anbindung der notwendigen Peripherie.

Zur physikalischen Integration in die bestehende Infrastruktur ist die Bridge mit dedizierten Schnittstellen ausgestattet. 
Für den DMX-Empfang wird ein RS-485-Transceiver (SN75176) verwendet, der als Pegelwandler zwischen den 
differenziellen DMX-Signalen nach RS-485-Standard und den TTL-Signalen des ESP32 agiert. 
Für den Art-Net-Betriebsmodus ist ein Ethernet-Modul vom Typ USR-ES1 integriert, welches den physikalischen 
Netzwerkzugriff für den Art-Net-Empfang gewährleistet.

Die Weiterleitung der verarbeiteten Steuerdaten an die Leuchtröhren erfolgt über das integrierte 
Wi-Fi-Modul des ESP32. Zur Erhöhung der Reichweite und Kommunikationsstabilität im störanfälligen Bühnenumfeld ist 
eine externe Antenne vorgesehen.

Die primäre Spannungsversorgung der Bridge-Komponenten erfolgt per 5V-USB-C-Anschluss direkt über das ESP32-Development-Board. 
Zur Erfüllung der Anforderungen bezüglich Systemkonfiguration und Benutzerfreundlichkeit verfügt das Bedienpanel über 
ein monochromes OLED-Display sowie drei Taster für die Menüführung. 
Es wurde exakt dieselbe Hardwarekombination wie bei der Röhre gewählt, um die Handhabung 
intuitiv und komponentenübergreifend einheitlich zu gestalten.

\subsection{LED-Leuchtröhren}
\label{subsec:led-leuchtroehren-hardware}

Die LED-Leuchtröhren unterliegen baulichen Beschränkungen aufgrund des schlanken Formfaktors bei einem Außendurchmesser von 40\,mm. 
Aus diesem Grund wurde für die Steuerung der \textit{XIAO ESP32-S3} Mikrocontroller ausgewählt. 
Dieser zeichnet sich durch einen besonders kleinen Formfaktor bei ausreichend verfügbaren GPIO-Pins 
und einem dedizierten Anschluss für eine externe WLAN-Antenne aus.

Als Leuchtmittel werden LED-Streifen des Typs WS2815 mit einer Dichte von 60 LEDs pro Meter verwendet. 
Der WS2815 teilt dasselbe Kommunikationsprotokoll wie der weitverbreitete WS2812B (5V-TTL-Logikpegel), 
verfügt jedoch über einen integrierten Spannungsregler und wird mit einer Arbeitsspannung von 12V betrieben. 
Die Anhebung der Betriebsspannung verringert den Spannungsabfall bei längeren Kabel- und
Streckenführungen. Zudem müssen die verwendeten Komponenten nur auf geringere
Stromstärken ausgelegt werden.

Der zusätzliche schaltungstechnische Aufwand zur Erzeugung der Betriebsspannung von 5V (für die restliche Elektronik) wird 
durch die Vorteile der 12V-LED-Architektur kompensiert. 
Zur Realisierung kommt ein kompakter DC-DC-Step-Down-Wandler zum Einsatz, der die von einem externen 12V-Netzteil 
bezogene Eingangsspannung auf 5V reduziert. 
Da der XIAO ESP32-S3 über einen eigenen On-Board-Spannungsregler verfügt, kann das 5V-Niveau direkt 
an dessen VCC-Eingangspin anliegen. 
Dieser transformiert die anliegende Spannung weiter auf die nominelle Chipspannung von 3,3V.

Für die Kommunikation kommen bei den Röhren die reziproken Gegenstücke der Bridge-Sender zum Einsatz: 
Das drahtlose Signal wird über die Wi-Fi-Schnittstelle des ESP32-S3 (mit externer Antenne) demoduliert.

Wie beim Steuergerät umfasst das Human-Machine-Interface auch hier ein monochromes OLED-Display 
sowie drei Navigationsschalter. 
Die konkreten Pinzuteilungen, Schaltungsdetails und Platinenentwürfe der Baugruppen
werden in Kapitel~\ref{chap:implementierung} erläutert.

\section{Software-Architektur}
\label{sec:software-architektur}

Das Software-Design des Systems folgt einem modularen und objektorientierten Ansatz. 
Ein zentraler Aspekt der Architektur ist die Verwendung derselben Codebasis für beide 
Geräteklassen (Bridge und Leuchtröhre). Durch bedingte Kompilierung wird mittels Umgebungsvariablen je nach 
Zielsystem entschieden, welche Logikblöcke eingebunden werden, was Redundanzen vermeidet und 
die Wartbarkeit erhöht.
Abbildung~\ref{fig:software-architektur} zeigt die resultierende Gliederung.

\begin{figure}[H]
    \centering
    \begin{tikzpicture}[x=1cm, y=1cm]
        \node[device, text width=4.7cm, align=center] (bfw) at (0,0)
            {\textbf{Bridge-Firmware}\\[2pt]\footnotesize DMX- und Art-Net-Eingang,
             Effektgenerierung im Master-Modus, Adressvergabe};
        \node[device, text width=4.7cm, align=center] (tfw) at (6.4,0)
            {\textbf{Leuchtröhren-Firmware}\\[2pt]\footnotesize Funkempfang, Auswertung
             der eigenen Adresse, Rendering};

        \node[modul] (m1) at (-1.3,-2.9) {Menü und Display};
        \node[modul] (m2) at (3.2,-2.9)  {Drahtlose Kommunikation};
        \node[modul] (m3) at (7.7,-2.9)  {LED-Steuerung};

        \begin{scope}[on background layer]
            \node[huelle, fit=(m1)(m2)(m3),
                  label={[font=\small\bfseries]below:Gemeinsame Module}] (hull) {};
        \end{scope}

        \draw[link] (bfw.south) -- (hull.north -| bfw.south);
        \draw[link] (tfw.south) -- (hull.north -| tfw.south);
    \end{tikzpicture}
    \caption{Gliederung der Software: gemeinsame Hardwareabstraktion, zur Übersetzungszeit ausgewählte gerätespezifische Firmware.}
    \label{fig:software-architektur}
\end{figure}

\subsection{Modulare Struktur und geteilte Komponenten}
\label{subsec:gemeinsame-softwarekomponenten}

Um die Kernfunktionen abzubilden, ist die Software in funktional getrennte Module unterteilt, 
die geräteübergreifend grundlegende Aufgaben zur Hardwareabstraktion übernehmen:

\begin{itemize}
    \item \textbf{Benutzerinteraktion (Menü und Display):} Kapselt die Verwaltung lokaler 
    Einstellungen, sichert Konfigurationen persistent und visualisiert Systemparameter für den Nutzer.
    \item \textbf{Drahtlose Kommunikation:} Stellt eine generische Abstraktionsschicht 
    für das Senden und Empfangen der systeminternen Datenpakete über das Funknetzwerk bereit.
    \item \textbf{Lichterzeugung (LED-Controller):} Übernimmt die logische Übersetzung von 
    rohen Datenwerten (Farben, Parameter oder Pixelarrays) in dynamische Lichtsignale
\end{itemize}

Ein übergeordneter Steuerungsprozess orchestriert diese Module, aktualisiert das System und 
etabliert – abhängig vom aktuellen Betriebsmodus – den direkten Datenfluss zwischen Peripherie 
und Aktoren.

\subsection{Gerätespezifische Firmware-Profile}
\label{subsec:geraetespezifische-firmware}

Aufbauend auf den gemeinsamen Modulen fokussieren sich Bridge und Leuchtröhre jeweils auf 
ihre dedizierten Aufgaben innerhalb der Systemtopologie:

\textbf{Bridge-Firmware:} \\
Die Software der Bridge agiert primär als Gateway. Sie initiiert 
zusätzliche spezialisierte Komponenten zum Einlesen und Verarbeiten externer Protokolle 
(DMX und Art-Net). Um Latenzen im Netzwerkverkehr zu minimieren, trennt das Architekturdesign 
Protokoll-Parsing vom regulären Ausgabezyklus. 
Dies ermöglicht das parallele Verarbeiten von eingehenden Datenströmen, während die Ausgabe 
an die Röhren kontinuierlich und ohne Unterbrechungen erfolgt.
Anschließend werden die aufbereiteten Datenpakete für die drahtlose Verteilung
bereitgestellt.

\textbf{Leuchtröhren-Firmware:} \\
Das Profil der Leuchtröhren zeichnet sich durch einen reaktiven, schlanken Ausführungszyklus 
aus, um Framedrops zu vermeiden und Latenzen zu minimieren:

\begin{itemize}
    \item \textbf{Autarker Betrieb:} Die Lichtsignalerzeugung agiert komplett losgelöst und speist 
    aktuelle User-Inputs in die eigene Lichtdarstellung ein.
    \item \textbf{Drahtloser Empfang:} Das System fragt fortlaufend den aktuellen Daten-Payload 
    für die eigene Geräteadresse aus der Funkschnittstelle ab und projiziert identifizierte Werte 
    umgehend auf die Pixel.
\end{itemize}

\section{Übertragungstechnologien und Protokolle}
\label{sec:uebertragungstechnologien}

Während Abschnitt~\ref{sec:datenfluss-und-signalverarbeitung} den inneren Signalablauf 
der Softwaremodule beschreibt, definiert dieser Abschnitt die netzwerktechnischen Konzepte und 
verwendeten Protokolle, um die Brücke zwischen den physikalischen Schnittstellen und 
den Software-Logikblöcken zu schließen.

\subsection{Drahtlose Übertragung über ESP-NOW}
\label{subsec:drahtlose-uebertragung}

Für die drahtlose Kommunikation wurde das herstellerspezifische
\textit{ESP-NOW}-Protokoll von Espressif gewählt. Seine technischen Eigenschaften sind in
Abschnitt~\ref{subsec:espnow-protokoll} beschrieben. Ausschlaggebend waren zwei
Eigenschaften, die unmittelbar auf die Anforderungen an Latenz und Stabilität einzahlen.

Erstens arbeitet das Protokoll verbindungslos. Damit entfallen nicht nur die
zeitaufwändigen Assoziierungsvorgänge eines Access-Point-Client-Modells, sondern auch die
Verbindungszustände, die im laufenden Betrieb verloren gehen und wiederhergestellt werden
müssten. Das System benötigt keine Netzwerkinfrastruktur und kann folglich auch nicht
durch deren Ausfall beeinträchtigt werden.

Zweitens adressiert die Bridge sämtliche Röhren über eine gemeinsame Broadcast-Adresse.
Der Sendeaufwand ist dadurch unabhängig von der Anzahl der Empfänger, und alle Röhren
erhalten denselben Datenstand innerhalb desselben Übertragungsfensters, was die
Voraussetzung für die in Abschnitt~\ref{subsec:synchronisation-roehem} beschriebene
Synchronisation ohne eigenes Zeitprotokoll. Der Broadcast-Betrieb kennt allerdings keine
Empfangsbestätigung. Die fehlende Zustellgarantie wird deshalb dadurch ausgeglichen, dass
die Bridge denselben Lichtzustand fortlaufend erneut aussendet: Geht eine Aussendung
verloren, trägt die nächste denselben Wert. Wie wirksam dieses Verfahren ist, untersucht
Kapitel~\ref{chap:evaluierung}.

\subsection{Kabelgebundene Schnittstellen}
\label{subsec:kabelgebundene-uebertragung}

Die Anbindung an die vorhandene Infrastruktur, Lichtpulte und Steuerungssoftware
erfolgt über DMX512 als primäres Eingangsprotokoll
(Abschnitt~\ref{subsec:dmx512}). Die Architektur nutzt dabei die Robustheit der
differentiellen RS-485-Übertragung, um die Signale störungssicher einzulesen, bevor sie
in das interne Format überführt werden. Damit fügt sich die Bridge ohne Anpassungen in
bestehende DMX-Strecken ein und erscheint dem Lichtpult als gewöhnliches Endgerät.

Alternativ nimmt die Bridge dieselben Steuerdaten über Art-Net entgegen
(Abschnitt~\ref{subsec:artnet}). Art-Net dient dabei ausschließlich als
Transportweg für ein DMX-Universum über eine Ethernet-Verbindung. Die Daten sind
inhaltlich identisch zu den über RS-485 empfangenen und durchlaufen ab der
Normalisierung denselben Verarbeitungspfad. Da die Bridge jedoch stets nur eine der beiden
Schnittstellen auswertet, ist die Wahl der Quelle als eigener Eintrag der Moduswahl
ausgeführt. Am Verhalten der Leuchtröhren und am weiteren Verarbeitungspfad ändert sich
dadurch nichts.

\subsection{Protokoll-Abstraktion (Internes Datenformat)}
\label{subsec:internes-protokoll}

Um die unterschiedlichen externen Eingabemethoden (lokales UI, DMX, Art-Net) zu harmonisieren, überführt die Systemarchitektur 
jegliche Eingaben stets in normierte interne Datenstrukturen (Data Transfer Objects), noch bevor sie verarbeitet oder weitergeleitet werden. 
Das Konzept zielt darauf ab, die physikalische Übertragungsebene von der darstellenden LED-Logik vollständig zu entkoppeln.
Durch die Kapselung in einen Transport-Frame, welcher stets eine Metainformation über die Startadresse mitführt, 
lassen sich DMX-Universen gezielt in kleinere Segmente fragmentieren. 
Dadurch muss nicht das gesamte Universum über die Bandbreiten-limitierte Funkschnittstelle übertragen werden, 
sondern lediglich jener Payload-Sektor, den die adressierten Röhren tatsächlich verarbeiten.

\section{Datenfluss und Signalverarbeitung}
\label{sec:datenfluss-und-signalverarbeitung}

Die Verarbeitung der Steuerdaten folgt unterschiedlichen Pfaden je nach Betriebsmodus. 
In diesem Abschnitt werden die End-to-End-Datenpfade, die Latenzcharakteristiken und die 
Synchronisationsmechanismen beschrieben.

\subsection{End-to-End-Datenpfade}
\label{subsec:end-to-end-datenpfade}

Der Datenfluss im System lässt sich entsprechend der Betriebsmodi in unterschiedliche Kanäle unterteilen:

\subsubsection{Bridge -- Master-Modus}

Im Master-Modus erzeugt die Bridge den Lichtzustand selbst. Aus den am Bedienmenü
gewählten Parametern entsteht fortlaufend ein Steuerdatensatz, der in das interne
Transportformat gekapselt und per Broadcast an alle Röhren verteilt wird. Die Röhren
werten den für sie bestimmten Abschnitt aus und setzen ihn unmittelbar in Licht um. Ein
externes Steuersignal ist an keiner Stelle des Pfades beteiligt.

\subsubsection{Bridge -- DMX-Modus}

Im DMX-Modus fungiert die Bridge als Gateway zwischen DMX-Strecke und Funkverbindung. Das
eingehende Universum wird über die RS-485-Schnittstelle empfangen, vollständig gepuffert
und anschließend auf denjenigen Adressbereich reduziert, den die angesteuerten Röhren
tatsächlich auswerten. Die Größe des übertragenen Pakets wächst dabei linear mit der
Anzahl der Röhren, da je Röhre vier Steuerbytes vorgesehen sind. Jede Röhre entnimmt dem
Broadcast den ihrer Startadresse zugeordneten Abschnitt.

\subsubsection{Leuchtröhre -- Standalone-Modus}

Ohne Verbindung zur Bridge erzeugt die Röhre den Lichtzustand aus den lokal gewählten
Parametern. Der Datenpfad verläuft vollständig innerhalb des Geräts, von der
Menüeingabe über die Lichterzeugung bis zur Ansteuerung des LED-Strips.

\subsubsection{Leuchtröhre -- DMX-Modus}

Im DMX-Modus ist der Empfang bewusst vom Ausgabezyklus entkoppelt: Eingehende Funkpakete
werden asynchron entgegengenommen und zwischengespeichert, während die Lichterzeugung
den jeweils aktuellen Pufferinhalt zyklisch ausliest. Diese Trennung stellt sicher, dass
der Empfang die Ausgabe nicht blockiert und umgekehrt keine Pakete verloren gehen, die
zwischen zwei Ausgabezyklen eintreffen.

\subsection{Datenpaketformate}
\label{subsec:datenpaketformate}

Auf Architekturebene werden zwei Datenkonzepte unterschieden, deren konkrete
speichertechnische Repräsentation Kapitel~\ref{chap:implementierung} beschreibt. Der
Steuerungsdatensatz fasst die Parameter einer einzelnen Röhre zusammen und ist damit die
kleinste adressierbare Einheit. Er wird im Master-Modus wie im DMX-Modus verwendet. Das
Transportformat kapselt eine Folge solcher Datensätze für die drahtlose Übertragung und
führt dabei die Startadresse des übertragenen Ausschnitts mit. Erst diese
Metainformation erlaubt es, statt des vollständigen Universums nur den tatsächlich
benötigten Ausschnitt zu übertragen.

\subsection{Latenzcharakteristiken}
\label{subsec:latenzcharakteristiken}

Die Systemlatenz von der Signaleingabe zur Lichtausgabe wird durch die Summierung 
folgender Teilschritte definiert:

\subsubsection{Bridge-Seite}

\begin{itemize}
    \item \textbf{Signal-Empfang und Decodierung:} Zeitaufwand für das physikalische Einlesen 
    des externen Protokolls (abhängig von der Bus-Baudrate des DMX- oder Netzwerk-Art-Net-Signals).
    \item \textbf{Pufferung und Datenformatierung:} Interne Verarbeitungszeit zur Konvertierung in das Systemprotokoll.
    \item \textbf{Funkübertragung:} Die Ausbreitungs- und Übertragungszeit über den drahtlosen Kanal zur Röhre.
\end{itemize}

\subsubsection{Röhren-Seite}

\begin{itemize}
    \item \textbf{Funksignal-Empfang:} Zeit für die Signalerkennung und Bereitstellung im Speicher.
    \item \textbf{Daten-Abfrage:} Dauer bis zur zyklischen Verarbeitung der bereitgestellten Daten durch die Lichterzeugungslogik.
    \item \textbf{LED-Refresh:} Dauer zur physikalischen Ansteuerung der LED-Hardware.
\end{itemize}

Das Systemarchitektur-Design stellt sicher, dass die Anforderungen an reaktive Beleuchtung 
(Abschnitt~\ref{subsec:echtzeitverhalten}) prinzipiell erfüllt werden können. 
Konkrete Implementierungsdetails und gemessene Zeiten sind Kapitel~\ref{chap:implementierung} 
und \ref{chap:evaluierung} zu entnehmen.

\subsection{Pufferstrategie und Datenaktualisierung}
\label{subsec:pufferstrategie}

Für die drahtlose Übertragung gilt eine Overwrite-Strategie nach dem Prinzip
\textit{Latest Data Only}: Ein eintreffendes Paket überschreibt den vorhandenen
Pufferinhalt unmittelbar, ohne dass ältere Stände aufbewahrt werden. Für Steuerdaten
einer Beleuchtung ist dies die angemessene Wahl, denn ein überholter Lichtzustand hat
keinen Wert mehr, sobald ein neuerer vorliegt. Die Entscheidung hat zwei Folgen, die
unmittelbar auf die Anforderungen einzahlen: Zu jedem Ausgabezyklus liegt stets der
aktuellste verfügbare Wert an, und ein verlorenes Paket führt nicht zu einem Ausfall,
sondern lediglich dazu, dass der zuletzt empfangene Zustand länger gehalten wird. Das in
Abschnitt~\ref{subsec:zuverlaessigkeit} geforderte definierte Verhalten bei
Verbindungsverlust ergibt sich damit unmittelbar aus der Pufferstrategie und erfordert
keine gesonderte Fehlerbehandlung.

Auf der Eingangsseite der Bridge werden dagegen vollständige Universen gepuffert, da ein
DMX-Frame erst nach seinem Ende vollständig ausgewertet werden kann. Eingehende Frames
werden strukturell geprüft, bevor der für die Röhren relevante Adressbereich entnommen
wird. Der Empfang über Art-Net ist vom Ausgabezyklus entkoppelt: Netzwerkpakete werden
asynchron entgegengenommen und zwischengespeichert, sodass sie dem Ausgabetakt ohne
zusätzlichen Jitter zur Verfügung stehen.

\subsection{Synchronisation mehrerer Röhren}
\label{subsec:synchronisation-roehem}

Um eine gleichzeitige Lichtausgabe bei verteilten Röhren zu gewährleisten, kombiniert 
die Architektur folgende Ansätze, anstatt auf aufwändige Zeitsynchronisationsprotokolle zurückzugreifen:

\begin{itemize}
    \item \textbf{Broadcast-Verteilung:} Dieser Ansatz übernimmt den Grundgedanken der DMX-Topologie.
    Das Leitsignal wird simultan an alle Empfangsknoten gesendet, 
    sodass die logische Anweisung an allen Röhren innerhalb desselben Übertragungsfensters am Receiver anliegt.
    
    \item \textbf{Logische Adressierungsschemata:} Die Broadcast-Frames bündeln die Informationen für mehrere Geräte. 
    Jedes Endgerät extrahiert über eine feste Adresszuweisung nur den für sich relevanten Datenblock.
    
    \item \textbf{Dezentrale Autonomie:} Lokal arbeiten die Geräte asynchron mit eigenen Refresh-Zyklen. 
    Die Ausführungsintervalle sind hierbei jedoch so kurz gewählt, dass visuelle Laufzeitunterschiede 
    zwischen Knotenpunkten weit unterhalb der physiologischen menschlichen Wahrnehmungsschwelle für Bewegungen (24 Hz also ca. 42 ms) bleiben \cite{fujii2018,ramunae2025}.
\end{itemize}

Diese Architektur ermöglicht eine dezentrale, eventgesteuerte Verarbeitung ohne zusätzliche Synchronisationsprotokolle 
und behält damit Einfachheit und Zuverlässigkeit bei.

\section{Konfiguration und Bedienkonzept}
\label{sec:konfiguration-und-bedienkonzept}

Bridge und Leuchtröhren werden ausschließlich am Gerät konfiguriert, ohne zusätzliche
Software und ohne Netzwerkzugang. Beide Gerätetypen verwenden dafür dasselbe
Bedienkonzept und dieselbe Kombination aus monochromem OLED-Display und drei Tastern,
sodass sich die Bedienung nicht unterscheidet, gleich welches Gerät der Anwender vor sich
hat. Die Navigation ist zweistufig aufgebaut: Die erste Ebene führt durch die verfügbaren
Parameter, die zweite dient deren Bearbeitung. Die konkrete Umsetzung einschließlich
Eingabeerfassung, Anzeige und Speicherung beschreibt
Abschnitt~\ref{subsec:menu}.

Architektonisch bedeutsam ist, dass der Parametersatz nicht fest vorgegeben ist, sondern
vom aktiven Betriebsmodus abhängt. Angezeigt werden ausschließlich die im jeweiligen
Modus wirksamen Einstellungen, was die Bedienung auf das Notwendige reduziert. Die
Moduswahl selbst ist davon ausgenommen und bleibt stets erreichbar. Daraus ergeben sich
die folgenden Parametersätze:

\begin{itemize}
    \item \textbf{Standalone-Modus (Leuchtröhre):} Effektauswahl und Farbwahl.
    \item \textbf{Master-Modus (Bridge):} Effektauswahl, Farbwahl, Anzahl der
          angesteuerten Leuchtröhren sowie Funkkanal.
    \item \textbf{DMX-Modus:} DMX-Startadresse der Leuchtröhre sowie Funkkanal.
    \item \textbf{Art-Net-Betrieb:} zusätzlich die Netzwerkkonfiguration der Bridge
          (IP-Adresse, DHCP-Aktivierung sowie Art-Net-Netz, Subnetz und Universum).
\end{itemize}

Sämtliche Werte werden persistent im nichtflüchtigen Speicher des Mikrocontrollers
abgelegt und beim Systemstart automatisch geladen. Die Geräte sind damit nach einer
Spannungsunterbrechung ohne erneute Konfiguration betriebsbereit, eine Anforderung, die
sich unmittelbar aus dem Veranstaltungsbetrieb ergibt
(Abschnitt~\ref{subsec:benutzerfreundlichkeit}).
